\documentclass[11pt]{article}

\usepackage[margin=1in]{geometry}
\usepackage{graphicx}
\usepackage{xcolor}
\usepackage{listings}
\usepackage{hyperref}
\usepackage{xurl}
\usepackage[htt]{hyphenat}
\usepackage{enumitem}
\usepackage{caption}
\usepackage{parskip}

\setlength{\emergencystretch}{3em}
\hypersetup{colorlinks=true, linkcolor=blue!50!black,
  urlcolor=blue!50!black, citecolor=blue!50!black}

\definecolor{codebg}{gray}{0.96}
\definecolor{promptbg}{rgb}{0.95,0.97,1.0}
\lstdefinestyle{code}{basicstyle=\ttfamily\small, backgroundcolor=\color{codebg},
  breaklines=true, frame=single, framerule=0pt, columns=fullflexible,
  showstringspaces=false, xleftmargin=4pt, xrightmargin=4pt}
\lstdefinestyle{prompt}{basicstyle=\ttfamily\small\itshape,
  backgroundcolor=\color{promptbg}, breaklines=true, frame=leftline,
  framerule=2pt, rulecolor=\color{blue!40}, columns=fullflexible,
  xleftmargin=8pt, aboveskip=6pt, belowskip=6pt}
\lstset{style=code}
\lstnewenvironment{prompt}{\lstset{style=prompt}}{}
\newcommand{\fig}[3]{\begin{figure}[ht]\centering
  \includegraphics[width=#2\linewidth]{figures/#1}
  \caption{#3}\label{fig:#1}\end{figure}}

\title{\textbf{Reconstructing Historic Fossil Caves in X3D from\\
Century-Old Survey Data: A Human--AI Case Study}}
\author{Alexander Hoffman \and Claude (Anthropic Claude Code)\\[2pt]
  \small\textit{Draft --- collaborative authorship}}
\date{\today}

\begin{document}
\maketitle

\begin{abstract}
We reconstruct two Pleistocene fossil caves of the McCloud River, California ---
\emph{Potter Creek Cave} and \emph{Samwel Cave}, both excavated 1903--06 under the
direction of the vertebrate palaeontologist John~C.~Merriam --- as interactive,
to-scale X3D scenes built procedurally from the original survey literature
(Sinclair 1904; Furlong 1906; Feranec et al.\ 2007). As in our companion
HAnim study, we treat the human--AI pair as co-authors: the human supplies
intent, domain pointers, family knowledge, and aesthetic and ethical judgment;
the AI supplies procedural generation, a tight verify-by-render loop, and
contributions back to the toolchain. We report three results. First, a method for
turning narrative cave surveys into standards-conformant X3D, organised around a
\emph{provenance discipline} that explicitly separates documented geometry from
interpretive dressing. Second, a physically-based cave look-development pass ---
displaced and normal-mapped limestone, layered ``god-ray'' light shafts, and
rippled water --- rendered in the X\_ITE web player. Third, and most useful to the
toolchain, an \textbf{X\_ITE rendering backend} for the X3D Model Context Protocol
(MCP) server that lets the tool \emph{see} HAnim humanoids and X3D~4.0
PhysicalMaterial (PBR) scenes that the previous X3DOM path could not render at
all under headless software WebGL. A personal and ethical thread runs throughout:
the excavator was the human author's great-great-grandfather, and the caves are
Winnemem Wintu sacred sites now drowned by Shasta Dam.
\end{abstract}

\section{Introduction}

This paper is a worked example of human--AI pairing applied to a problem that is
equal parts engineering, archival scholarship, and cultural care. Over two
sessions a human collaborator and an AI agent (Anthropic's Claude, in Claude
Code) built two interactive X3D scenes that reconstruct the chambers of
\emph{Potter Creek Cave} and \emph{Samwel Cave} in Shasta County, California ---
to the dimensions recorded in their original early-twentieth-century surveys, and
viewable in a web browser.

The motivation was personal. The caves were excavated between 1903 and 1906 by
field parties under \textbf{John Campbell Merriam} (1869--1945), the University of
California vertebrate palaeontologist who later led the Carnegie Institution and
helped found Save the Redwoods --- and the human author's great-great-grandfather.
The reports that resulted (Sinclair 1903, 1904, 1905; Furlong 1906) are among the
founding documents of West Coast Pleistocene palaeontology. We set out to ask: can
a human--AI pair turn those century-old prose surveys into faithful, inspectable
3D, honestly labelled as to what is measured versus imagined, using only the
open Web3D toolchain?

The contribution is less the scenes than the \emph{method and the discipline}:
every artifact was produced through the X3D MCP server, validated against the X3D
schema and a set of semantic checks, and visually verified by the agent rendering
each scene and reading the resulting image. Along the way we found and fixed the
single largest gap in that loop --- the MCP could not actually render the modern
material model --- and contributed the fix upstream.

\fig{pcc_hero}{0.86}{Potter Creek Cave, hero view near the entrance third.
Documented to Sinclair (1904): the 107\,ft chamber, $\sim$75\,ft domed roof, the
two coalescing breccia fans, and the 42\,ft entrance pit (left, with its daylight
shaft). Interpretive: speleothems, pools, rock texture, and lighting.}

\section{The sites and their sources}

\paragraph{Potter Creek Cave.} Sinclair's 1904 monograph gives a chamber
``one hundred and seven feet long, about thirty feet wide at its widest part,
with the roof rising about seventy-five feet above the lowest point of the
floor.'' Two ``fan-like deposits of earth and stalagmite-cemented breccia''
slope from opposite ends and coalesce in the middle, with near-vertical chimneys
above each fan apex; access required ``a vertical descent of forty-two feet'' by
rope ladder. The galleries trend northwest--southeast, parallel to the strike of
the McCloud limestone, at 1500\,ft elevation. The deposit yielded roughly 52
vertebrate species, 21 of them extinct --- including the type specimens of the
shrub-ox \emph{Euceratherium collinum} (a new genus and species named from this
cave), abundant remains of the giant short-faced bear \emph{Arctodus simus}, the
Shasta ground sloth \emph{Nothrotheriops shastensis}, \emph{Megalonyx},
\emph{Canis dirus}, mammoth, camel, and horse.

\paragraph{Samwel Cave.} Five kilometres north on the same McCloud arm, Samwel
Cave (Wintu \emph{sawal}, ``holy place'') is a branching two-level system in the
McCloud Limestone at 460\,m elevation. Furlong's 1906 survey, redrawn in Feranec
et al.\ (2007, Fig.~2), names its spaces: the Entrance and Porcupine Entrance, the
large Pleistocene Hall, the Gate, Chamber One, Merriam's Chamber, and Chamber Two
(``Furlong's Room''), reached only by a deep drop. Tradition records the cave as
the \emph{Cave of the Lost Maiden} --- the Wintu girl Olchanolmet, who fell to her
death down a ``90-foot-deep hole'' --- and as the \emph{Cave of the Magic Pools},
where ``Wintu medicine men bathed in the water pools to get magic strength.''
Nearly 1000 specimens (45 mammal, 13 bird species) span the Last Glacial Maximum
($\sim$17{,}100--23{,}600 cal~BC).

\paragraph{An ethical note.} Both caves are Winnemem Wintu sacred sites. Shasta
Dam (1945) flooded the McCloud River canyon --- villages, burial and ceremonial
places --- and the caves now sit on the Shasta Lake shore. We name the Winnemem
connection, the Lost Maiden, and the Magic Pools explicitly and with attribution,
as context rather than scenery, and we surface the dam-flooding in every scene's
caption.

\section{Method: from narrative survey to standards-conformant X3D}

Each cave is generated procedurally by a single Python program emitting X3D~4.0
XML (\texttt{generate\_cave.py}, \texttt{generate\_samwel.py}). Three ideas did
most of the work.

\paragraph{Sectioning.} A closed cave shell is a dark blob from outside and is
unlit from within. We instead build the \emph{back half} of each surface (the
$z\le0$ portion), opening the chamber toward the viewer --- the same longitudinal
and cross-section conventions the 1904 and 1906 reports used to draw the caves on
paper. This makes the interior simultaneously legible and lightable. Potter Creek
becomes a lofted longitudinal section; Samwel becomes a labelled cross-section of
ellipsoidal chambers connected by passages, echoing Furlong's plate.

\paragraph{Physically-based look-development.} Walls and floor use X3D~4.0
\texttt{PhysicalMaterial} with a tileable metallic-roughness texture set ---
albedo, tangent-space normal, and metallic-roughness maps generated from
value-noise with NumPy/PIL (\texttt{generate\_cave\_textures.py}) --- combined
with multi-octave geometric displacement of the mesh for true silhouette relief.
On top sit procedural speleothems (stalactites, stalagmites, fused columns,
flowstone draperies), breakdown blocks, rippled PBR pools, and layered emissive
``god-ray'' cones with drifting dust motes down the open shafts.

\paragraph{Provenance discipline.} Because the geometry is part measured and part
imagined, every generator carries a \texttt{PROVENANCE} block, and every scene a
caption, that state plainly which features are \emph{to scale from the survey}
(chamber dimensions, the fans, the chimneys, the 42\,ft pit; Samwel's chamber
topology, the $\sim$90\,ft hole, the Magic Pools) and which are
\emph{interpretive} (all speleothems, the rock texture, pool extents, lighting,
and --- for Samwel --- the individual room sizes, which the sources do not give).
We regard this explicit separation as a small but reusable contribution to honest
heritage and scientific reconstruction.

\fig{samwel_section}{0.9}{Samwel Cave as a labelled cross-section after Furlong
(1906) / Feranec et al.\ (2007). Documented topology: Entrance and Porcupine
Entrance, Pleistocene Hall, the Gate, Chamber One, Merriam's Chamber, and Chamber
Two down the $\sim$90\,ft hole to the sacred Magic Pool. Room sizes and surface
detail are interpretive.}

\section{The render loop and the X3D MCP server}

Our governing practice is \emph{verify by render}: a scene that passes both the
XSD schema check and our semantic checks can still be visually wrong --- off
camera, unlit, mis-scaled, or silently dropped. The canonical loop is therefore
\texttt{describe\_node} (look up exact fields and containerFields) $\rightarrow$
build $\rightarrow$ \texttt{validate\_x3d} \emph{and} \texttt{validate\_semantic}
$\rightarrow$ \textbf{render and look}. Both cave scenes pass both validators
cleanly.

The decisive finding of today's session was that the loop's last step was broken
for exactly the scenes we care about. The MCP's \texttt{render\_image} tool used
X3DOM with the synchronous Playwright API. Two failures compounded: (i) called
from an asynchronous MCP tool it raised \emph{``Sync API inside the asyncio
loop''}; and (ii) X3DOM under headless software WebGL (SwiftShader) clears the
background but \emph{draws no geometry at all}, and in any browser it does not
implement X3D~4.0 PhysicalMaterial or HAnim. The tool literally could not see a
PBR cave or an articulated humanoid.

We rebuilt the backend around \textbf{X\_ITE}, which renders both. The new
implementation (i) uses asynchronous Playwright so it runs inside the MCP event
loop; (ii) serves the scene over an ephemeral loopback HTTP server, because
X\_ITE fetches the \texttt{.x3d} by XHR and browsers block that over
\texttt{file://}; and (iii) when given a file path, serves that file's own
directory so relative assets (PBR textures, inlined geometry) resolve as they do
in a browser. With this in place the agent could finally inspect its own output;
every figure in this paper was produced through it.

\fig{pcc_section}{0.92}{Potter Creek Cave, full longitudinal section, rendered
through the new X\_ITE MCP backend. Visible: displaced, normal-mapped limestone;
the two breccia fans; the three shafts; columns and dripstone; the banded
stratigraphic column; and the daylight shaft down the great pit.}

\section{What we learned today}

\begin{enumerate}[leftmargin=1.4em]
\item \textbf{Renderer choice is a correctness property, not an aesthetic one.}
For X3D~4.0 PBR and HAnim, X3DOM is not merely lower-fidelity --- under our
headless setup it renders \emph{nothing}, and it omits the material model
entirely. X\_ITE is the faithful player. An authoring tool that cannot render the
standard's current material model cannot verify its own work.
\item \textbf{Transport matters.} \texttt{file://} silently defeats X\_ITE (CORS);
inline X3D inside HTML is mangled by the HTML parser's lowercasing. A tiny local
HTTP server is the reliable path, and pages should \emph{detect} \texttt{file://}
and say so rather than show a black screen.
\item \textbf{Look up, do not recall.} Every material field
(\texttt{baseTexture}, \texttt{normalTexture}, \texttt{metallicRoughnessTexture};
\texttt{metallic} whose default is $1$ and must be set to $0$ for rock) came from
\texttt{describe\_node}, not memory. Guessing containerFields is the most common
silent failure, and the validators that name the exact fix are worth more than
the schema check.
\item \textbf{Provenance is cheap and clarifying.} Writing down what is measured
versus imagined took minutes and removed a whole class of future
mis-citation --- and it forced better research (it is how we caught that the
maiden's fall was a documented \emph{90-foot} drop, and corrected the pit from a
shallow notch to a true shaft connecting the two levels).
\item \textbf{Get the legend right.} A late prompt --- ``make sure you read the
map legends correctly'' --- sent us back to the published cross-section to verify
chamber identities and to discover the famous Magic Pool, which then became the
emotional and geometric focus of the Samwel descent view.
\end{enumerate}

\fig{samwel_hero}{0.86}{Samwel Cave, the descent: Chamber One, the $\sim$90\,ft
hole the Lost Maiden fell down (the only connection between the upper and lower
levels), and Chamber Two with the glowing Magic Pool.}

\section{Ideas and future work}

\begin{itemize}[leftmargin=1.4em]
\item \textbf{Populate the caves with their fauna.} The whole point of these
caves was what came out of them. The procedural-creature pipeline already built
for a rigged moose could place the actually-excavated megafauna --- the type
\emph{Euceratherium}, the short-faced bear, the Shasta ground sloth --- at their
recorded depths, turning each cave into a labelled palaeontological diorama.
\item \textbf{Ground-truth from the field.} A planned in-person visit to Samwel
could supply photographs of the real chambers and formations, letting us replace
estimated room sizes and invented dripstone with measured shapes --- closing the
interpretive gap the provenance block currently flags.
\item \textbf{Stratigraphic interactivity.} Expose the documented strata
(Potter Creek's clay/ash/breccia; Samwel's flowstone--breccia--gravel column) as
a scrubbable cut, with fauna pinned to levels and radiocarbon ages on hover.
\item \textbf{Toward true volumetrics.} Our god-rays and ``refractive'' water are
faked with layered translucent geometry; X3D lacks participating-media
scattering. A volumetric or screen-space approach in X\_ITE would be a genuine
rendering contribution.
\item \textbf{A reusable survey$\rightarrow$X3D recipe.} The section-cut + PBR +
provenance pattern generalises to any cave or excavation with a published plan;
it could become a small documented template for the AI-with-X3D community.
\end{itemize}

\section{Contributions back to the toolchain}

Beyond the scenes, the session produced concrete improvements to the Web3D X3D
MCP server, carried on a dedicated branch and covered by tests: the X\_ITE
rendering backend described above (async, loopback-served, asset-aware);
semantic validators for containerField (with the canonical field suggested),
USE-before-DEF ordering, and interpolator key/keyValue length; an
\texttt{autofix\_x3d} tool that rewrites containerFields and returns the corrected
document; granular-mode parity for non-default node placement; and proactive
server \texttt{instructions} that frame the canonical describe$\rightarrow$build
$\rightarrow$validate-both$\rightarrow$render path for fresh model instances. We
also maintain an upstream \texttt{x3d.py} bug report documenting a dropped
\texttt{containerField} on non-default fields and a non-spec
\texttt{EnvironmentLight.global} default --- both of which serialise to
schema-valid XML that nonetheless renders incorrectly.

\section{Conclusion}

Two caves that a family forebear opened a hundred and twenty years ago now exist
as honest, inspectable, standards-conformant 3D, made entirely through the open
Web3D toolchain by a human and a model working as a pair. The engineering lesson
is blunt: an authoring agent is only as good as its ability to see what it made,
and for modern X3D that means X\_ITE. The human lessons --- provenance,
attribution, and care for whose sacred ground this is --- were supplied by the
person, and they are the part that matters most.

\paragraph{Reproducibility.} All generators, textures, scenes, viewer pages, and
the full render history are in the project repository; \texttt{./start\_caves.sh}
serves the launcher locally. Documented-versus-interpretive provenance is recorded
in each generator and on screen.

\end{document}
